% =====================================================================
% main-bilingual.tex -- LINGSHU 中英对照版（XeLaTeX）
% 结构：每个 section 先英文、后中文；中文文件在 sections-zh/，不含 \label。
% 构建：xelatex main-bilingual && bibtex main-bilingual && xelatex ×2
% source of truth 仍是 main.tex（英文）；本文件只做对照阅读，不改内容。
% =====================================================================
\documentclass[11pt]{article}

\usepackage[a4paper,margin=0.9in]{geometry}
\usepackage{fontspec}
\usepackage{xeCJK}
\setmainfont{Times New Roman}
\setCJKmainfont[Path=C:/Windows/Fonts/, Extension=.ttf]{simhei}
\setCJKmonofont[Path=C:/Windows/Fonts/, Extension=.ttf]{simhei}
\usepackage{amsmath}
\usepackage{amssymb}
\usepackage{booktabs}
\usepackage{graphicx}
\usepackage{microtype}
\usepackage[numbers,sort&compress]{natbib}
\usepackage[colorlinks=true,linkcolor=black,citecolor=blue,urlcolor=blue]{hyperref}
\usepackage{url}
\usepackage{seqsplit}
\usepackage{listings}
\lstset{basicstyle=\ttfamily\footnotesize,breaklines=true,columns=fullflexible,
  keepspaces=true,frame=none,aboveskip=4pt,belowskip=4pt}

\setlength{\parskip}{2pt}
\frenchspacing

\newcommand{\system}{\textsc{LingShu}}
\newcommand{\host}{\textsc{Dsh}}
\newcommand{\eg}{e.g.,\ }
\newcommand{\ie}{i.e.,\ }
\newcommand{\wrt}{w.r.t.\ }

% 每节：英文在前，中文在后；中文缺失时降级为提示，不中断编译。
\newcommand{\pair}[2]{%
  \input{#1}%
  \IfFileExists{#2}{\input{#2}}{%
    \par\medskip\noindent\textit{（本节中文对照待补：该节由并行工作流新增。）}\par}}
\newcommand{\zhlang}[1]{\par\medskip\noindent{\large\bfseries #1}\par\nopagebreak}

\title{\system: A Cybernetic Control Plane for Auditable\\
Cognition Plugins in Agent Hosts\\[6pt]
\large 灵枢：面向 Agent 宿主可审计认知插件的控制论控制面\\[2pt]
\normalsize\itshape 中英对照版（Bilingual draft）}
\author{\system{} Project Team}
\date{October 2026}

\begin{document}
\maketitle

\noindent\textbf{Reading note.} Each section appears in English first and in
Chinese immediately after. Numbers, qualifiers, negations and boundary
statements are translated one-for-one; the English source
(\texttt{main.tex}) remains authoritative if the two ever diverge. Sections
added after the first translation pass are marked when their Chinese text is not
yet available.

\zhlang{阅读说明}
每节先英文后中文。数字、限定词、否定句与边界声明逐条对应；若中英出现歧义，以英文源文件（\texttt{main.tex}）为准。首次翻译之后新增的章节若暂无中文，会在原位置标注。

% ---------------- Abstract ----------------
\begin{abstract}
Persistent AI companions are where many users meet long-context agents: someone
raises a digital character over hundreds of turns, then it forgets its name or
catchphrase, and closing the session leaves it gone. The failure is less often
capability than the injection surface: memories are read as current facts,
stylistic signals soften refusals, persona constraints decay, and truncation can
push safety instructions out. We present \system{}, a deterministic and auditable
control plane that treats prompt assembly as a control surface. Six plugins share
one read-only constitutional boundary; a five-domain composer arbitrates among
them; and the control plane adds integral feedback with anti-windup, phase-shift
preemptive monitoring with persona anchors, requisite-variety reduction, and an
STPA-derived gate that structurally constrains what reaches the model at the
assembly boundary. Exactly one node is fail-closed---the STPA audit---so a
violation blocks injection instead of degrading it; every other component
degrades open and records why. The contract is machine-checkable: 892 Python
assertions across 14 files and 396 plugin-side assertions, with cross-copy
parameter consistency checked key by key. A chaos suite exercises each mechanism
in simulated and real hosts. We call the union the \emph{cyber-cognitive
envelope}; the name is ours, not a new mechanism.
\end{abstract}

\begingroup
\renewcommand{\abstractname}{摘要}
\IfFileExists{sections-zh/01-abstract.tex}{\begin{abstract}
持久化的 AI 伴侣（persistent AI companions）是许多用户接触长上下文 Agent 的入口：
有人用数百轮对话养一个数字角色，然后它忘了自己的名字或口头禅，关掉会话它就不在了。
这种失效较少来自能力，更多来自注入面（injection surface）：记忆被当作当前事实来读，
风格信号软化了拒绝，人格约束逐渐衰减，而截断可能把安全指令挤出去。
我们提出 LINGSHU——一个把 prompt 组装视为控制面的确定性、可审计控制面。
六个插件共享同一个只读的宪法边界；一个五域合成器在它们之间裁决；
控制面在模型调用之前加入带抗饱和（anti-windup）的积分反馈、带人格锚点的相变前馈监控、
必要多样性降维，以及一道源自 STPA、在组装边界上结构性约束「什么能到达模型」的闸门。
恰好一个节点是 fail-closed——STPA 审计——因此违规会阻断注入而不是降级注入；
其余组件以 fail-open 方式降级并记录原因。契约可机器校验：14 个文件中的 892 条 Python 断言
与 396 条插件侧断言，副本之间的参数一致性按键比对。
一个混沌套件在模拟宿主与真实宿主上逐一驱动每个机制；v5.4 的同回合投递路径让相位状态
落在携带临界比率的那次请求里，并由 24 条独立回归断言钉住。一次真实宿主 smoke 运行显示
两层已经解耦：闸门以 soft 重定向拦下一次角色扮演越狱，模型则在人设内拒绝，而不是复述系统规则。
我们把这一并集称为\emph{赛博认知包络}（cyber-cognitive envelope）；
名字是我们起的，不是一种新机制。
\end{abstract}
}{}
\endgroup

% ---------------- 1 Introduction ... 8 Conclusion ----------------
\pair{sections/02-intro}{sections-zh/02-intro}
\pair{sections/03-related}{sections-zh/03-related}
\pair{sections/04-problem}{sections-zh/04-problem}
\pair{sections/05-method}{sections-zh/05-method}
\pair{sections/06-experiments}{sections-zh/06-experiments}
\pair{sections/06b-supplement}{sections-zh/06b-supplement}
\pair{sections/06c-host}{sections-zh/06c-host}
\pair{sections/06d-engineering}{sections-zh/06d-engineering}
\pair{sections/06e-actuator}{sections-zh/06e-actuator}
\pair{sections/06f-baselines}{sections-zh/06f-baselines}
\pair{sections/06g-case-study}{sections-zh/06g-case-study}
\pair{sections/06h-immersive-case}{sections-zh/06h-immersive-case}
\pair{sections/07-discussion}{sections-zh/07-discussion}
\pair{sections/08-limitations}{sections-zh/08-limitations}
\pair{sections/09-conclusion}{sections-zh/09-conclusion}

% ---------------- Appendix ----------------
\appendix
\pair{sections/10-appendix}{sections-zh/10-appendix}
\pair{sections/10b-appendix-audit}{sections-zh/10b-appendix-audit}
\pair{sections/10c-appendix-tool}{sections-zh/10c-appendix-tool}
\pair{sections/10d-appendix-hostpilot}{sections-zh/10d-appendix-hostpilot}
\pair{sections/10e-risk-classifier}{sections-zh/10e-risk-classifier}
\pair{sections/10f-appendix-supplychain}{sections-zh/10f-appendix-supplychain}

\bibliographystyle{plainnat}
\bibliography{refs}

\end{document}
